\section{Background and Motivation}
\label{sec:background}

\subsection{CPS Communication and Attack Surface}

Industrial PLCs execute control logic at 1--10\,ms scan cycles, reading sensors, computing outputs, and updating actuators on every cycle. Communication with SCADA/HMI systems occurs over field protocols (S7, Modbus, OPC-UA) that expose process data blocks containing all tag values. On many deployed systems---including the Siemens S7-1200 in default configuration---these data blocks are readable and writable by any network-reachable host without authentication.

This creates a well-documented attack surface: an adversary with S7-level access can read the live tag stream (sensor values, actuator states, setpoints) and inject writes at the process-semantic level. Real-world CPS attacks from Stuxnet~\cite{langner2011stuxnet} to PIPEDREAM~\cite{cisa2022pipedream} exploited exactly this---manipulating process data within existing protocols rather than deploying novel exploits. Stuxnet wrote to PLC data blocks~\cite{falliere2011w32}; Industroyer issued standard IEC~60870-5-104 commands~\cite{cherepanov2017industroyer}; the Oldsmar attacker changed a single setpoint~\cite{cisa2021oldsmar}. The trend is toward reusable, protocol-aware attack frameworks: PIPEDREAM provides modular OT-specific tools targeting multiple PLC families~\cite{greenberg2022pipedream}.

\subsection{LLMs for CPS Security: State and Gap}

LLMs have been applied to CPS security in several ways: AttackLLM~\cite{ahmed2025attackllm} generates process-aware attacks on static traces; LLM4PLC~\cite{fakih2024llm4plc} synthesizes PLC invariants; LLM-based IDS approaches~\cite{ann2024ids_llm,sharshar2025explainable} classify historical traffic logs. In the adjacent domain of autonomous cyber agents, PentestGPT~\cite{deng2024pentestgpt} and Hackphyr~\cite{rigaki2024hackphyr} automate IT penetration testing, and MALF~\cite{ning2025malf} fuzzes industrial protocols at the packet layer.

However, none of these approaches embed an LLM in a \emph{live} PLC control loop---observing real tag values, reasoning about the process phase in real time, and executing writes whose physical consequences are immediately observable. Without closed-loop execution it remains unclear whether LLM-generated strategies survive practical constraints (timing, process physics, controller overwrite), and whether defenses improve meaningfully after exposure to such attacks.

\subsection{PLC-in-the-Loop Evaluation}

CPSForge addresses this gap by combining a real Siemens S7-1200 PLC with Factory~I/O, a software-based process emulator. Every tag value the LLM observes is a real S7 protocol read; every write is a real S7 write that changes the controller's data block. Unlike offline evaluations, the LLM sees what a real adversary would see: live tag values updating every 500\,ms, with real controller dynamics. This \emph{PLC-in-the-loop} design is essential for studying online-vs-static attack modes (RQ3), LLM defense under real timing (RQ4), and whether rich context changes model behavior on live systems (RQ1).